%%
%% Test: Mathematical Environments — v1.1
%% Stage: math refactor (stages 1-4)
%% 
%% Purpose: Verify math settings:
%%   - Equation numbering (Iranian standard: formula-chapter)
%%   - Math operators (grad, curl, diver, rank, ...)
%%   - Number sets (\N, \Z, \Q, \R, \C)
%%   - Delimiters (\abs, \norm, \ceil, \floor, \inner, \paren)
%%   - Sets (\set, \setbuilder, \given)
%%   - Matrices (\mat, \pmat, \vmat)
%%   - Derivatives (\dd, \deriv, \pderiv)
%% 

\documentclass{matinbook}

\title{تست محیط‌های ریاضی — نسخه ۱.۱}
\author{تیم MatinBook}
\date{۱۴۰۵}

\booktitle{تست ریاضیات}

\begin{document}

\maketitle

\chapter{بررسی محیط‌های ریاضی}
\label{chap:test-math}

این سند، محیط‌های ریاضی \texttt{mb-math} را بررسی می‌کند.

%----------------------------------------
\section{شماره‌گذاری معادلات}
\label{sec:equation-numbering}

معادله‌ی اول در فصل اول:

\begin{equation}
    \label{eq:math-test-1}
    f(x) = x^2 + 2x + 1
\end{equation}

معادله‌ی دوم در فصل اول:

\begin{equation}
    \label{eq:math-test-2}
    g(x) = \int_{0}^{\infty} e^{-t^2} \, dt = \frac{\sqrt{\pi}}{2}
\end{equation}

\textbf{بررسی:} شماره‌ی معادلات باید به فرمت \texttt{(فرمول-فصل)} باشند، مثلاً \texttt{(۱-۱)} و \texttt{(۲-۱)}.

%----------------------------------------
\section{عملگرهای ریاضی}
\label{sec:operators}

\textbf{حساب برداری:}
\[
    \grad f, \quad \curl \vec{F}, \quad \diver \vec{F}
\]

\textbf{جبر خطی:}
\[
    \rank(A), \quad \diag(a_1, a_2, \ldots, a_n), \quad \trace(A)
\]

\textbf{بهینه‌سازی:}
\[
    \argmax_{x \in \R} f(x), \quad \argmin_{x \in \R} g(x)
\]

\textbf{نظریه اعداد:}
\[
    \lcm(a, b), \quad \gcdop(a, b)
\]

\textbf{توابع خاص:}
\[
    \sinc(x), \quad \sgn(x), \quad \erf(x)
\]

%----------------------------------------
\section{مجموعه‌های اعداد}
\label{sec:number-sets}

\[
    \N \subset \Z \subset \Q \subset \R \subset \C
\]

%----------------------------------------
\section{دلimiters}
\label{sec:delimiters}

\[
    \abs{x}, \quad \norm{\vec{v}}, \quad \ceil{x}, \quad \floor{x}, \quad \inner{u}{v}, \quad \paren{a + b}
\]

%----------------------------------------
\section{مجموعه‌ها}
\label{sec:sets}

\[
    \set{1, 2, 3}, \quad \setbuilder{x \in \R}{x > 0}
\]

%----------------------------------------
\section{ماتریس‌ها}
\label{sec:matrices}

\[
    \mat{1 & 2 \\ 3 & 4}, \quad \pmat{a & b \\ c & d}, \quad \vmat{1 & 2 \\ 3 & 4}
\]

%----------------------------------------
\section{مشتقات}
\label{sec:derivatives}

\[
    \deriv{f}{x}, \quad \pderiv{f}{x}, \quad \dd{x}
\]

%----------------------------------------
\section{محیط‌های چندخطی}
\label{sec:align}

\begin{align}
    a &= b + c \\
    d &= e + f \\
    g &= h + i
\end{align}

%----------------------------------------
\section{نتیجه‌گیری}
\label{sec:conclusion}

اگر این سند بدون خطا کامپایل شود:

\begin{itemize}
    \item ✅ شماره‌گذاری معادلات به فرمت \texttt{(فرمول-فصل)} است
    \item ✅ عملگرهای ریاضی کار می‌کنند
    \item ✅ مجموعه‌های اعداد کار می‌کنند
    \item ✅ دلimiters کار می‌کنند
    \item ✅ ماتریس‌ها کار می‌کنند
    \item ✅ مشتقات کار می‌کنند
    \item ✅ محیط \texttt{align} کار می‌کند
\end{itemize}

\textbf{وضعیت تست:} قبول

\end{document}
